\documentclass[options]{philosophersimprint}
%\usepackage{sectsty}
%\linespread{1.2}
\usepackage{xunicode}
\usepackage{xltxtra}
\defaultfontfeatures{Mapping=tex-text}
\setmainfont[
Ligatures={Common}, Numbers={OldStyle}, Variant=01,
BoldFont=LinLibertine_RB.otf,
ItalicFont=LinLibertine_RI.otf,
BoldItalicFont=LinLibertine_RBI.otf
]{LinLibertine_R.otf}
\setmonofont[Scale=0.8]{DejaVuSansMono.ttf}
\usepackage{stmaryrd}
\usepackage[toc,page]{appendix}
\usepackage{amsmath,amsthm,amssymb}
\usepackage[colorlinks=true, citecolor=blue, linkcolor=blue, hyperfootnotes=false]{hyperref}
\usepackage{nameref}
\usepackage{soul}
\usepackage{frame}
\usepackage[style=apa, natbib=true]{biblatex}
\usepackage{scrextend}
\usepackage{mathrsfs}
%\addtokomafont{labelinglabel}{\sffamily}
\makeatletter
\newcommand\labelitem[2]{%
  \item[#1]\def\@currentlabel{#1}\label{#2}}
\makeatother
\usepackage[inline]{enumitem}
\usepackage{framed, csquotes}
\usepackage{tikz}
\usetikzlibrary{cd}
\tikzset{%
    symbol/.style={%
        draw=none,
        every to/.append style={%
            edge node={node [sloped, allow upside down, auto=false]{$#1$}}}
    }
}
%\usepackage{anyfontsize}
\usepackage[nodayofweek,level]{datetime}
\usepackage{etoolbox}
\patchcmd{\formatdate}{,}{}{}{}
\usepackage{enumitem, hyperref}
\makeatletter
\def\namedlabel#1#2{\begingroup
    #2%
    \def\@currentlabel{#2}%
    \phantomsection\label{#1}\endgroup
}
\makeatother
%\usepackage{anyfontsize}
% 0.95in rather than the usual 1in: this handout runs right at the four-page
% boundary, and the extra strip is what holds it there.
\usepackage[margin=0.95in]{geometry}
\usepackage{pifont}% http://ctan.org/pkg/pifont
\newcommand{\cmark}{\ding{51}}%
\newcommand{\xmark}{\ding{55}}
\usepackage{forest}
\forestset{
  L1/.style={draw=black,},
  L2/.style={,edge={,line width=0.8pt}},
}
\usepackage[acronym]{glossaries}
\usepackage{stackengine}
\stackMath
\usepackage{scalerel}
\usepackage{titlesec}
\newcommand\halo{{{\circ}\mkern-1mu}}
\usepackage{linguex}
\usepackage{parskip}
\usepackage{fancyhdr}
\pagestyle{plain}
% Diagrams
%\usepackage{pgfplots}
\usepackage{tikz}
\usetikzlibrary{arrows,calc,patterns,positioning,shapes}
\usetikzlibrary{decorations.pathmorphing}
\tikzset{
modal/.style={>=stealth’,shorten >=1pt,shorten <=1pt,auto,
node distance=1.5cm,semithick},
world/.style={circle,draw,minimum size=1cm,fill=gray!15},
point/.style={circle,draw,fill=black,inner sep=0.5mm},
reflexive/.style={->,in=120,out=60,loop,looseness=#1},
reflexive/.default={5},
reflexive point/.style={->,in=135,out=45,loop,looseness=#1},
reflexive point/.default={25},
}
\usepackage{tree-dvips}
\usepackage{qtree}

\DeclareMathOperator*{\argmax}{arg\,max}
\DeclareMathOperator*{\argmin}{arg\,min}
\newtheorem{theorem}{Theorem}
\newtheorem{corollary}{Corollary}
\newtheorem{lemma}{Lemma}
\newtheorem{proposition}{Proposition}
\theoremstyle{definition}
\newtheorem{definition}{Definition}
\newtheorem{example}{Example}
\newtheorem{fact}{Fact}
\newtheorem{remark}{Remark}
\newtheorem{claim}{Claim}
\newcommand\numberthis{\addtocounter{equation}{1}\tag{\theequation}}

\newcommand{\ev}[2][]{\ensuremath{\llbracket #2\rrbracket^{#1}}}
%command
\newcommand{\A}{\mathscr{A}}
\newcommand{\cred}{\mathcal{C}}
\newcommand{\des}{\mathcal{D}}
\newcommand{\scred}{\mathsf{C}}
\newcommand{\sdes}{\mathsf{D}}
\newcommand{\ut}{\mathcal{U}}
\newcommand{\rb}{{\rrbracket}}
\newcommand{\lb}{{\llbracket}}
\newcommand{\counterfactual}{\ensuremath{%
  \Box\kern-1.5pt
  \raise1pt\hbox{$\mathord{\rightarrow}$}}}
  \newcommand{\tr}{\mathsf{T}}
%\renewcommand{\fa}{\mathsf{F}}


\usepackage{FOL_Yale/notation}

\begin{document}

PHIL 1115: First-order Logic \hfill Yale, Fall 2026\\
Lecture 2, Handout \hfill Dr. Daniel Grimmer

\textbf{0. Review: Truth Preserving by Virtue of its Form}

Last lecture I motivated one (among many) possible definitions of a ``good argument''. To a formal logician, an argument is good (valid) if it is truth preserving by virtue of its form.
\begin{align}
\text{Either $p_1$ or $q_1$,\qquad not $p_1$\qquad}
&\text{$\therefore$ $q_1$}\\
\text{If $p_2$ then $q_2$,\qquad\qquad  $p_2$\qquad}
&\text{$\therefore$ $q_2$}\\
\text{If $p_3$ then $q_3$,\qquad  not $q_3$\qquad}
&\text{$\therefore$ not $p_3$}
\end{align}
Valid argument forms are like little machines: If we specify the non-logical terms (i.e., the $p$'s and $q$'s) such that the premises are all true, then the argument's formal validity guarantees that making the same specifications in the conclusion will make it true as well.

\textit{Formal Validity, read forwards}: if truth goes in, then truth comes out.

The very same guarantee can be read backwards. If a \textit{false} conclusion ever comes out, then a false premise must have gone in (for had all the premises been true, the conclusion would have been true too).

\textit{Formal Validity, read backwards}: if a falsity comes out, then a falsity went in.

Importantly, however, the machine only makes promises when it is fed good input. If garbage goes in (i.e., a false premise), then we have no guarantee about the truth/falsity of the conclusion. Garbage might come out (i.e., a false conclusion) although, as a matter of luck, a truth might come out too.

Logicians only care about building and understanding these little ``truth-preserving machines''. It is up to the user (e.g., biologists, politicians, your uncle) to come up with good premises (e.g., biological, political and uncle-ical truths). Hence, we have a division of labor between humans. Interestingly, we also seem to have a division of labor within the words of a valid argument.

\textbf{1. Identifying the Logical Vocabulary of English}

Recall our valid (but unsound) argument from the previous lecture:
\begin{itemize}
    \item[] All lizards are mammals.
    \item[] No mammals are egg-layers.
    \item[$\therefore$] No lizards are egg-layers.
\end{itemize}
Consider what the words ``lizards'', ``mammals'', and ``egg-layers'' mean and what they point to or correspond to out in the world. Importantly, these pieces of non-logical vocabulary play no role whatsoever in establishing the formal validity of the argument. Validity comes from the argument's logical form:
\begin{itemize}
    \item[] All Xs are Ys. \qquad No Ys are Zs. \qquad $\therefore$ No Xs are Zs.
\end{itemize}
Here then is the division of labor between the two kinds of words.
\begin{itemize}
    \item[$\bullet$] The job of the logical vocabulary (e.g., ``all'', ``are'', and ``no'') is to give the argument a valid form, i.e., to make it a good little truth-preserving machine. In this particular argument, they succeed. The argument is valid.
    \item[$\bullet$] The job of the non-logical vocabulary (e.g., ``lizards'', ``mammals'', and ``egg-layers'') is to make the premises true. In this particular argument, they fail. The argument is unsound.
\end{itemize}

In \textit{formal} logic, we are only interested in valid arguments: those which are truth preserving by virtue of their form. Importantly, the \textit{formal} validity of an argument rides entirely on the arrangement of its \textbf{logical vocabulary}: words like ``if\ldots then'', ``and'', ``or'', ``not'', ``all'', and ``no'', etc.

\textbf{Challenge Question} But if we limit our attention to \textit{formal} validity, what other notions of validity are we missing out on? Can't the \textit{content} of the non-logical vocabulary sometimes contribute to the validity of an argument? Is the division of labor always so strict as I have made it out to be? Try to find an argument which is valid (where the truth of the premises \textit{guarantees} the truth of the conclusion) for reasons other than its logical form. Besides logical structure, what else can provide this sort of guarantee?

Proceeding with formal validity: What exactly is the division between logical and non-logical vocabulary? Here is a provocative idea: \textit{the logical vocabulary of English consists of the domain-general ``scaffolding'' which appears uniformly across all subject matters.} Consider the words: ``all'', ``are'', ``no'', ``if ... then'', ``or'', etc. You are just as likely to find these words in a newspaper article about the stock market, as you are in a peer-reviewed biology journal, or a political opinion piece. This is not true of the words: ``lizards'', ``mammals'', and ``egg-layers''. These words are domain-specific.

So perhaps the logical/non-logical divide is exactly the same as the domain-general/specific divide. Another point of evidence in this direction is the existence of \textit{cross-domain propositions}:\footnote{This is closely related to the famous Frege-Geach Problem in metaethics, which suggests that moral claims must possess the same underlying logical semantics as descriptive claims in order to function correctly inside of conditionals or conjunctions.} ``Murder is wrong, and the sky is blue''.

\textbf{Challenge Question}. The above discussion suggests \textbf{Logical Monism:} the claim that there is \textit{One True Logic} that governs all correct reasoning, within and between all subject matters. This would position logic as a uniquely foundational topic which sits between and underneath all other subjects. But is this actually true? Do the exact same rules of logic apply perfectly across every conceivable domain? The system we will study in this course (First-Order Logic) is often treated as this universal standard. But one must remain skeptical: many philosophers deny Logical Monism. Moreover, many who accept it deny that First-Order Logic is the One True Logic. You should form your own opinion on this throughout the course.

\textbf{2. Picking Logical Vocabulary for our New Language} To study logical form with full precision, we now build a self-contained formal language whose logical vocabulary is explicitly fixed. The very first decision we face is \textit{which} logical terms to include, and this is a genuine choice, not something forced upon us.

For propositional logic we will include exactly five logical symbols: The \textbf{connectives}: $\Neg$, $\ConjTight$, $\Disj$, $\Cond$, and $\Bicond$.
\begin{itemize}
    \item $\Neg$ is \textbf{negation} and stands (roughly) for `It is not the case that ...'
    \item $\ConjTight$ is \textbf{conjunction} and stands (roughly) for the English `and'.
    \item $\Disj$ is \textbf{disjunction} and stands (roughly) for the English `or' (inclusive).
    \item $\Cond$ is known as the \textbf{material conditional} and stands (\textit{very} roughly) for the English `If... then...'.\footnote{We'll usually abbreviate the last two, referring to them just as `the conditional' and `the biconditional', respectively.}
    \item $\Bicond$ is known as the \textbf{material bi-conditional} and stands (roughly) for the English `...if and only if...'.
\end{itemize}
All of these ``roughly'' qualifiers will be discussed in detail in Lecture 3.

This choice of symbols deliberately leaves a great deal out. The quantificational words ``all'' and ``no'' (the very words doing the work in our lizard argument above) are \textit{not} among our connectives; the logic of quantification is the business of the second half of this course. Words like ``possibly'' and ``necessarily'' are left out too: their logic, \textit{modal logic}, belongs to another course entirely.

However, settling on these five connectives (and postponing the rest) we can already begin to translate some sentences from English into our formal language. For example, `If it's raining outside, then the streets are wet' can be translated simply as $p \Cond q$. Similarly, `The moon is made of green cheese and my cat is the President' can be translated as $r \Conj s$.\footnote{Actually, officially, this isn't quite right. See the notes about parentheses below.}

English has several ways of expressing a conditional, and they do not all put the antecedent first: `$q$, if $p$', `$p$ only if $q$', `$p$ is sufficient for $q$', `$p$ is necessary for $q$'. Which of these point the arrow which way is a question about what $\Cond$ \textit{means}, so we sort the whole family out next lecture, once $\Cond$ has a truth table. It \textit{will} appear on the exams.

\textbf{Challenge Question}. We chose these five connectives, but why these five? Why not include new symbols for ``implies'', ``but'', ``is a sufficient condition for'', etc.? Presumably, the ``neither ... nor'' construction can be handled by some combination of $\Neg$ and $\Disj$ (meaning ``not'' and ``or''), but what about the others? Exactly how expressive are these five symbols? Are they all necessary? That is, could we be just as expressive with fewer symbols? As we shall see in the second problem set, we can express every truth functional relationship using only $\ConjTight$, $\Disj$ and $\Neg$.

\textbf{3. The Language of Propositional Logic (Syntax)} We start with a stock of symbols. You can think of these symbols as the alphabet of the language.
\begin{itemize}
    \item Lower-case Roman letters possibly with subscripts (e.g., $p, q_2, r_7$). These are the \textbf{atomic formulas} of the language. They are placeholders for pieces of non-logical vocabulary (namely atomic sentences, defined below).
    \item Left and right parentheses: $($ and $)$. Think of these as the language's punctuation marks. They will help us parse long complicated expressions unambiguously.
    \item The five connectives, $\Neg, \ConjTight, \Disj, \Cond, \Bicond$, introduced above.
\end{itemize}

Any finite string of the foregoing symbols is called an \textbf{expression} of the language. Here are some expressions: \quad $)p\Neg\Disj$ \quad and \quad $(pq\Disj s$. Most expressions are ill-formed (ungrammatical) and so are meaningless.

However, what we're interested in primarily are \textbf{(well-formed) formulas} of the language. We build formulas up from the foregoing symbols using a recursive procedure:
\begin{itemize}
    \item[(i)] Every lower case letter (potentially with a subscript) is a formula of the language. For instance, $p$ is one and $r_5$ is another.
    \item[(ii)] Let $A$ represent a generic expression of the language.\footnote{Note something important here: $A$ is \textit{not} itself an expression of our language (nor a formula). That's because our language doesn't contain \textit{upper-case} Roman letters; it only contains lower-case ones. Rather, $A$ is a part of the meta-language which I'm using to talk \textit{about} expressions of our language. We'll talk more about this idea later on.} If $A$ is a formula of the language, then so is $\Neg A$.
    \item[(iii)] Let $A$ and $B$ represent two generic expressions of the language. If $A$ and $B$ are formulas of the language, then: \begin{itemize}
        \item $(A \Conj B)$ and $(A \Disj B)$ and $(A \Cond B)$ and $(A \Bicond B)$ are each formulas of the language.
    \end{itemize}
    \item[(iv)] Nothing else is a formula of the language.
\end{itemize}

We can use (i-iii) to build up complicated formulas of our language as follows.
\begin{itemize}
    \item By (i) $p$ is a formula of our language. Hence, by (ii) so is $\Neg p$
    \item By (i) $q$ is a formula of our language. Combining this with $\Neg p$ we can use (iii) to see that $(q \ \Disj\Neg p)$ is a formula of our language.
    \item We can also use (iii) to build $(p \Conj q)$. And then we can combine these together as $((q \ \Disj\Neg p)\Cond (p \Conj q))$.
\end{itemize}

To demonstrate that some expression is a formula of our language, one simply has to show how it can be constructed step by step using (i-iii).

But how can one demonstrate that some expression \textit{is not a formula of our language}? Note that (ii) and (iii) are purely positive, they validate formulas using pre-existing formulas. And (i) provides us with a stock of building blocks. The purpose of clause (iv) is to provide some negativity: no other constructions are allowed. (This point will be discussed further in Lecture 5.)

To demonstrate that some expression is NOT a formula of our language, one has to show that there is \textit{no possible way} to construct it using (i-iii).

Exercise: Demonstrate that $(p \Cond (p \Cond p)$ is not a formula of our language.

Exercise: Demonstrate that $p \Cond (p \Cond p)$ is also not a formula of our language.

As we shall now discuss, \textit{unofficially}, we are allowed to drop the outermost parentheses, which makes this last expression okay. Let me stress that \textit{officially it is not a well-formed formula}.

\textbf{4. Parentheses in Unofficial Formulas and the Main Connective} For the sake of readability, however, we'll sometimes omit the outermost parentheses (officially, they're still there; they're just written in invisible ink).\footnote{The full rule is: you can unofficially drop the outermost parentheses \textit{so long as the main connective is not a negation}. Consider these two official formulas with clearly different meanings: $\Neg(p \Conj q)$ and $(\Neg p \Conj q)$. Dropping the ``outermost parenthesis'' would yield $\Neg p \Conj q$ in both cases. Hence we disallow it in the first case, when the $\Neg$ is the main connective.} So, for example, we'll write $p \Cond (q  \Conj r)$ in place of $(p \Cond (q\Conj r))$. Also, if a formula contains many parentheses, sometimes we'll cheat by using square brackets in place of some of them (they're still officially parentheses; they're just poorly drawn). Thus, we might write $s\Bicond[p \Cond (q \Conj r)]$ in place of $s\Bicond(p \Cond (q\Conj r))$.

Note that you're not allowed to drop parentheses willy-nilly, even unofficially. To see why, consider this sentence: `The butler and the maid or the gardener committed the murder' formalized as $b \Conj m\Disj g$. This could mean two things:
\begin{itemize}
    \item `The butler committed the murder and (the maid or the gardener committed the murder)'. Formalized as $b \Conj (m\Disj g)$.
    \item `(The butler and the maid committed the murder) or the gardener committed the murder'. Formalized as $(b \Conj m)\Disj g$.
\end{itemize}
Such ambiguities would arise in our new language if we dropped parentheses willy-nilly. So we'll only ever drop \textit{outermost} parentheses, never inner ones.

Notice that the paths to building $b \Conj (m\Disj g)$ and $(b \Conj m)\Disj g$ are substantially different. In the first case, the last connective to be added was $\ConjTight$ whereas in the second case the last connective to be added was $\Disj$.  This demonstrates the general fact that every non-atomic formula in our language has a \textbf{main connective}: the last connective added in its construction.

%\textbf{Challenge Question}. Given any formula of our language, how can we be sure that there are not two different ways to build it, ending with different kinds of steps. If this were possible then there could be two (or more) ways of picking out the formula's main connective. But this never happens. Why?

\textbf{5. Subsentences and Propositional Forms} We now have a formal language of our own. The final task of this lecture is to connect it back to English: given an English sentence, how do we uncover the \textbf{propositional form} that displays its logical structure? To do this systematically, we need one more notion.

First, some terminology: a \textbf{subsentence} of a sentence, $S$, is a part of $S$ (or the whole of $S$) which is itself a complete sentence. For instance, consider the sentence, `If it's raining outside, then the streets are wet'. Hopefully, two subsentences jump out at you here: `It's raining.' and `The streets are wet.'

Sometimes subsentences contain other subsentences. Consider $S =$`If the moon is made of green cheese, then either my cat is the President or it's always sunny in Philadelphia.' The sentence $S$ has five subsentences:
\begin{enumerate}
    \item `It's always sunny in Philadelphia.'
    \item `My cat is the President.'
    \item `Either my cat is the President or it's always sunny in Philadelphia.'
    \item `The moon is made of green cheese.'
    \item `If the moon is made of green cheese, then either my cat is the President or it's always sunny in Philadelphia.'
\end{enumerate}
That's right! The whole sentence counts as a subsentence of itself. Think of it as being like $\leq$, less-than-or-equals.

An \textbf{atomic sentence} is a sentence which has no subsentences (other than itself). An example of an atomic sentence would be `My cat is the President'. Three of $S$'s subsentences are atomic.

Sometimes we have to rephrase a sentence to get it to reveal its logical structure more clearly. Consider `I had a donut and a coffee'. Intuitively, we can rephrase this as `I had a donut and I had a coffee' and thereby arrive at the propositional form `$p$ and $q$'. WARNING: It is not trivial to decide what kinds of minor rephrasing are allowed. Extracting propositional forms from English sentences is a bit of an art. Consider `Dan and Caroline got married in Iceland'. Does this mean the same thing as `Dan got married in Iceland and Caroline got married in Iceland'? And a connective can hide with no connective-word at all: `A louder alarm would have woken Sam' means `if the alarm had been louder, then Sam would have woken', the subjunctive `would' doing an `if's work. Whenever you rephrase like this, say which rephrasing you have adopted: the propositional form you end up with depends on it.

Once we have a propositional form, we can use it to produce English sentences by replacing its schematic letters $p$, $q$, etc. with complete English sentences. This produces an \textbf{instance} of the propositional form. An \textbf{argument form} is a list of propositional forms (the premises), together with one further propositional form (the conclusion).

\end{document}
